% Generated draft: requires mathematical and source review.
\documentclass[11pt]{article}
\usepackage[margin=25mm]{geometry}
\usepackage{fontspec}
\setmainfont{Latin Modern Roman}
\usepackage{amsmath,amssymb,mathtools,unicode-math}
\setmathfont{Latin Modern Math}
\usepackage{enumitem,needspace,xurl,hyperref}
\setlength{\parindent}{0pt}
\setlength{\parskip}{5pt}
\setlength{\emergencystretch}{3em}
\newcommand{\N}{\mathbb{N}}
\newcommand{\Z}{\mathbb{Z}}
\newcommand{\Q}{\mathbb{Q}}
\newcommand{\R}{\mathbb{R}}
\newcommand{\C}{\mathbb{C}}
\newcommand{\F}{\mathbb{F}}
\newcommand{\A}{\mathbb{A}}
\newcommand{\PP}{\mathbb{P}}
\newcommand{\E}{\mathbb{E}}
\newcommand{\eps}{\varepsilon}
\newcommand{\dd}{\,\mathrm{d}}
\begin{document}
\section*{UnsolvedMath ID: 30001193}
\textbf{Source problem code:} catalogue entry
\subsection{Definitions and mathematical statement}
Let \(A\) be a noetherian ring, as in the supplied source. The source does not specify whether \(A\) is commutative or whether modules are left or right modules; the constructions below are relative to a fixed module side, with the corresponding noetherian hypothesis. These choices require clarification before fixing the scope of the problem.

Write \(D(A)\) for the unbounded derived category of the chosen category of \(A\)-modules. Two possible interpretations of the source notation are
\[
D_{\mathrm{fg}}(A)=\{X\in D(A):H^n(X)\text{ is finitely generated for every }n\in\Z\}
\]
and
\[
D^b_{\mathrm{fg}}(A)=\{X\in D_{\mathrm{fg}}(A):H^n(X)=0\text{ for all but finitely many }n\in\Z\}.
\]
The source calls its category \(D_f(A)\), describing finitely generated cohomology without explicitly requiring boundedness. Its reference to compact objects makes this omission material; neither interpretation is silently selected here.

Let \(K(\operatorname{Inj}A)\) denote the homotopy category of unbounded cochain complexes of injective \(A\)-modules, whose morphisms are chain maps modulo chain homotopy. A thick subcategory of a triangulated category is a full subcategory closed under isomorphisms, shifts, distinguished triangles, and direct summands. A localizing subcategory is a full triangulated subcategory closed under isomorphisms and arbitrary set-indexed coproducts. An object \(C\) in a triangulated category \(\mathcal T\) with coproducts is compact if, for every set-indexed family \((X_i)_{i\in I}\), the canonical map
\[
\bigoplus_{i\in I}\operatorname{Hom}_{\mathcal T}(C,X_i)
\longrightarrow
\operatorname{Hom}_{\mathcal T}\!\left(C,\bigoplus_{i\in I}X_i\right)
\]
is an isomorphism. Write \(\mathcal T^c\) for the full subcategory of compact objects. A localizing subcategory generated by compact objects means the smallest localizing subcategory containing some set of objects of \(\mathcal T^c\).

\textbf{Classification problem, with unresolved source qualifications.} For every noetherian ring \(A\) in the intended scope, classify the thick subcategories of the intended category \(D_f(A)\). The supplied record also asks to classify all localizing subcategories of \(K(\operatorname{Inj}A)\).

The record presents these requests as equivalent via compact objects. This equivalence requires qualification: an identification with compact objects does not by itself identify a classification of thick subcategories with a classification of all localizing subcategories. Localizing subcategories generated by ambient compact objects must be distinguished from arbitrary localizing subcategories. The record supplies neither a restriction to the former nor hypotheses ensuring that every localizing subcategory has that form. Thus the two classification requests are retained separately, and their asserted equivalence remains unresolved.
\subsection{Short English statement}
For a noetherian ring \(A\), classify the thick subcategories of the source's derived category \(D_f(A)\), and classify the localizing subcategories of the homotopy category of injective complexes. The module side, commutativity assumptions, cohomological boundedness, and claimed equivalence of these requests need clarification. The supplied status is \emph{partially solved}; this is a statement-only draft requiring human mathematical review.
\subsection{Sources}
\begin{itemize}
\item Supplied local UnsolvedMath registry record, permanent ID \(30001193\), exact source problem code \texttt{catalogue entry}, title \emph{Thick Subcategories of Noetherian Derived Categories}; collection: Oberwolfach Reports Open Problems. Local registry references supplied with the task are \texttt{data/unsolvedmath/problems.json} and the supplementary \texttt{catalog/problems.json}; their contents were not independently inspected here.
\item \emph{Mini-Workshop: Support Varieties} (2009), Oberwolfach Reports, pp.\ 549--592, DOI \texttt{10.4171/owr/2009/10}: \url{https://doi.org/10.4171/owr/2009/10}. This is also the supplied literature-source link labeled ``original Oberwolfach problem.'' The record attributes the problem to workshop open-problem discussions and gives proposed year 2009; no individual proposer is supplied.
\item Supplied literature reference labeled ``Benson--Iyengar--Krause'': \url{https://doi.org/10.24033/asens.2076}. No title, theorem number, or precise statement of applicability is supplied.
\item Literature assessment dated exactly \texttt{2026-08-21}: supplied dated curation, not independently verified. The record labels the problem \emph{partially solved} and states that classification is complete for regular rings and important special classes such as finite-group algebras and complete intersections, while no general classification of thick subcategories of \(D_f(A)\) or localizing subcategories of \(K(\operatorname{Inj}A)\) is available. The background describes this as literature triage derived from supplied Oberwolfach curation and explicitly requests independent verification of time-sensitive status and cited claims. The scope and hypotheses of these special-case assertions are not supplied.
\item The registry supplies the statement-status label \texttt{corrected\_verified} and says its mathematical content was checked against the official report, with notation, context, and question form normalized without changing the source-level claim. This is a supplied verification assertion, not an independent verification performed for this draft.
\end{itemize}
\end{document}
